\documentclass[11pt]{article}

\usepackage[a4paper,margin=1in]{geometry}
\usepackage{amsmath,amssymb,amsthm,mathtools}
\usepackage{graphicx}
\usepackage{microtype}
\usepackage[hidelinks]{hyperref}

\newtheorem{theorem}{Theorem}
\newtheorem{lemma}{Lemma}
\newtheorem{corollary}{Corollary}
\newtheorem{remark}{Remark}

\title{Compact Abelian Torsion Groups Have Bounded Exponent}
\author{Candidate full answer for arXiv:1306.1480\\
Run \texttt{fa\_banach\_001}, lane 7}
\date{12 August 2026}

\begin{document}
\maketitle

\begin{center}
\fbox{\parbox{0.91\textwidth}{\textbf{Status.}
Candidate full proof, likely valid, pending specialist human review.
Under the standard meaning of ``torsion group'', the bounded-exponent
hypothesis in the source theorem is automatic.}}
\end{center}

\section{Source question}

Czuron and Wojciechowski prove a quantitative nonisomorphism theorem for a
``commutative compact torsion group with finite exponent'' and then ask
whether the assumption of bounded exponent is significant
\cite[p.~3]{CzuronWojciechowski}.

\begin{figure}[ht]
  \centering
  \includegraphics[width=\textwidth]{figures/open_problem_crop.png}
  \caption{The open remark in arXiv:1306.1480, PDF page 3.}
\end{figure}

For the class literally stated in the theorem, the answer is no: compactness,
commutativity, and torsion already force one common exponent.

\section{The topological group lemma}

We write compact groups additively.  A group is \emph{torsion} if every one
of its elements has finite order, and it has \emph{finite exponent} if some
positive integer annihilates every element.

\begin{lemma}\label{lem:bounded-exponent}
Every compact Hausdorff abelian torsion group has finite exponent.
\end{lemma}

\begin{proof}
Let $G$ be such a group.  For each $n\geq1$, set
\[
  K_n=\ker(n\,{\rm id}_G)=\{x\in G:nx=0\}.
\]
Because multiplication by $n$ is a continuous homomorphism, $K_n$ is a
closed subgroup of $G$.  The torsion assumption gives the countable cover
\[
  G=\bigcup_{n=1}^{\infty}K_n.
\]
Every compact Hausdorff space is a Baire space.  Hence some $K_N$ has
nonempty interior.  A subgroup with nonempty interior is open: if $U$ is a
nonempty open subset of $K_N$ and $u\in U$, then $U-u$ is a neighbourhood of
$0$ contained in $K_N$, and all translates of that neighbourhood show that
$K_N$ is open.

The quotient $G/K_N$ is therefore both compact and discrete, so it is finite.
Let $q$ be its exponent (one may simply take $q=|G/K_N|$).  For every $x\in
G$, the coset $x+K_N$ is killed by $q$, which says $qx\in K_N$.  Applying $N$
now gives
\[
  Nq,x=0.
\]
Thus the single integer $Nq$ annihilates all of $G$.
\end{proof}

\begin{corollary}\label{cor:source}
In the main theorem of arXiv:1306.1480, the phrase ``with finite exponent''
can be deleted without changing the class of groups.  In particular, for
every compact abelian torsion group $G$ satisfying the paper's subgroup
exclusion, all of the theorem's distortion conclusions and quantitative
lower bounds hold exactly as stated.
\end{corollary}

\begin{proof}
Apply Lemma~\ref{lem:bounded-exponent} to $G$, then invoke the source theorem.
\end{proof}

\section{Why the hypotheses matter}

The Baire step is the entire point.  Compactness cannot simply be dropped:
the discrete abelian group
\[
  \bigoplus_{m=2}^{\infty}\mathbb Z/m\mathbb Z
\]
is torsion but has no finite exponent.  Commutativity is used to ensure that
$K_n$ is a subgroup; for the source theorem the group is explicitly
commutative.

There is also an important terminological boundary.  A profinite or
zero-dimensional compact group need not be torsion: for example,
$\prod_{m\ge2}\mathbb Z/m\mathbb Z$ contains elements of infinite order.
Thus the argument resolves the question under the standard algebraic
meaning of the words appearing in the source theorem.  If ``torsion'' was
intended nonstandardly to mean merely profinite, or if the intended object
was instead an unbounded-exponent discrete dual group, that would be a
different problem not formulated by the displayed theorem.

\section{Verification and review recommendation}

The proof uses only four exact facts:
\begin{enumerate}
  \item the torsion condition gives $G=\bigcup_n K_n$;
  \item each $K_n$ is a closed subgroup;
  \item Baire category makes one $K_N$ open;
  \item a compact discrete quotient is finite.
\end{enumerate}
No quantitative estimate, classification theorem, or unproved regularity
claim is needed.

Review as a likely valid full answer to the literal question.  The only
substantive review point is semantic: confirm that ``compact torsion group''
is being used in its standard sense (every element has finite order).  Under
that reading, the finite-exponent assumption is provably redundant.

\begin{thebibliography}{9}
\bibitem{CzuronWojciechowski}
A. Czuron and M. Wojciechowski,
\emph{On the isomorphisms of Fourier algebras of finite abelian groups},
arXiv:1306.1480, 2013.
\url{https://arxiv.org/abs/1306.1480}.
\end{thebibliography}

\end{document}
